\providecommand{\kBKM}{\kappa_{\mathrm{BKM}}}
\providecommand{\kch}{\kappa_{\mathrm{ch}}}
\providecommand{\kcat}{\kappa_{\mathrm{cat}}}
\providecommand{\ClaimStatusTheorem}{{\bfseries\sffamily[Theorem]}}
\providecommand{\ClaimStatusConjectured}{{\bfseries\sffamily[Conjectured]}}
\providecommand{\Sp}{\mathbf{Sp}}
\providecommand{\GammaT}[2]{\Gamma_{#1}(#2)}

\section*{Wave~10 T4/20 --- The eight Gritsenko--Cl\'ery dd-Siegel forms (full tuple table)}
\label{sec:wave10-t4-8-form-dd-siegel}

\paragraph{Slide source.}
Lorgat 2020 presentation deck (\texttt{/tmp/automorphic-slides.txt}),
slide ``Diagonal Divisor Siegel Modular Forms'' (lines~45--54), quoting
the Gritsenko--Cl\'ery theorem verbatim; §7--8 development reappears at
lines~856--953 (``Arithmetic Geometry of some CY3's'', Conjecture~1).

\paragraph{Slide theorem, verbatim.}
\ClaimStatusTheorem{}
``For congruence subgroups $\Gamma_t(N) < \Gamma_t$ there are exactly
eight dd-modular forms'' (slide line~47--48), displayed as:
\begin{center}
$M_5(\Gamma_1, \nu_2) = \Delta_5$ \quad
$M_2(\Gamma_2, \nu_4)$ \quad
$M_3(\GammaT{1}{2}, \nu_2)$
\\[2pt]
$M_1(\Gamma_3, \nu_6)$ \quad
$M_2(\GammaT{1}{3}, \nu_2)$
\\[2pt]
$M_{1/2}(\Gamma_4, \nu_8)$ \quad
$M_{3/2}(\GammaT{1}{4}, \nu_4)$ \quad
$M_1(\GammaT{2}{2}, \nu_4)$
\end{center}
Here $\Gamma_t$ denotes the integral paramodular group of polarisation
$t$ (slide lines~30--37), $\GammaT{t}{N}$ its Hecke-type congruence
subgroup of level~$N$, and $\nu_k$ a primitive character of order~$k$.

\paragraph{Complete $(t, N, k, \chi, c_k(0))$ tuple table.}
\ClaimStatusTheorem{} (primary source Gritsenko--Cl\'ery, Compos.\ Math.\
149 (2013), Theorem~1.1, reprised in the slide).
\begin{center}
\small
\begin{tabular}{lcccccl}
\hline
Form & $t$ & Level $N$ & Weight $k$ & Mult.\ $\chi$ & $c_k(0) = 2k$ & $\kBKM$ target \\
\hline
$M_5(\Gamma_1,\nu_2) = \Delta_5$      & $1$ & $1$ & $5$   & $\nu_2$ & $10$ & Vol~III $N{=}1$ \\
$M_2(\Gamma_2,\nu_4)$                 & $2$ & $1$ & $2$   & $\nu_4$ & $4$  & Vol~III $N{=}2$ \\
$M_3(\GammaT{1}{2},\nu_2)$            & $1$ & $2$ & $3$   & $\nu_2$ & $6$  & Vol~III $N{=}3$ \\
$M_1(\Gamma_3,\nu_6)$                 & $3$ & $1$ & $1$   & $\nu_6$ & $2$  & -- \\
$M_2(\GammaT{1}{3},\nu_2)$            & $1$ & $3$ & $2$   & $\nu_2$ & $4$  & Vol~III $N{=}4$ \\
$M_{1/2}(\Gamma_4,\nu_8)$             & $4$ & $1$ & $1/2$ & $\nu_8$ & $1$  & -- \\
$M_{3/2}(\GammaT{1}{4},\nu_4)$        & $1$ & $4$ & $3/2$ & $\nu_4$ & $3$  & Vol~III $N{=}6$ \\
$M_1(\GammaT{2}{2},\nu_4)$            & $2$ & $2$ & $1$   & $\nu_4$ & $2$  & -- \\
\hline
\end{tabular}
\end{center}

The first column reproduces the slide's form labels exactly (lines~50--54);
the $(t, N)$ split reads off the group subscript---$\Gamma_t$ has $N=1$,
$\GammaT{t}{N}$ has the displayed~$N$. The weight~$k$ is the integer or
half-integer subscript of~$M_{\cdot}$. The multiplier $\chi$ is the
character $\nu_{\cdot}$. The datum $c_k(0) = 2k$ records the Borcherds
constant coefficient by the universal Borcherds-weight identity (see
below).

\paragraph{Why exactly eight.}
\ClaimStatusTheorem{} (Gritsenko--Cl\'ery 2013, proof of Theorem~1.1
pp.~1567--1571). The Jacobi-to-Siegel additive-multiplicative lift sends
a weight-$0$ index-$t$ weak Jacobi form $\phi$ with pole of order
exactly~$1$ along the diagonal to a weight-$k$ dd-Siegel modular form
$M_k(\Gamma_t(N), \chi)$. The classification forces
$tN \leq 4$ (Wave-9 L4 lattice obstruction, matching the Gritsenko--Cl\'ery
proof's bound derived from the $\Gamma_t(N)$-translate orbit counting
of $\mathbb{H}_1$-diagonals in $\mathbb{H}_2$). Enumerating
$(t, N)$ with $tN \leq 4$ and $\gcd(t, N) \in \{1, 2\}$ produces exactly
eight pairs; each pair admits a unique dd-form up to scalar, with
$k = k(t, N)$ and $\chi = \chi(t, N)$ determined by the Jacobi-form
constant-coefficient constraint.

\paragraph{Grothendieck context.}
The eight forms are the $H^0$-generators of the ring of
dd-Siegel-cuspidal forms across the tower
$\{\Gamma_t(N)\}_{tN \leq 4}$; the $tN \leq 4$ bound is the paramodular
analogue of the Grothendieck finiteness theorem for sections of a line
bundle on a projective variety with ample divisor class---the
dd-polar class $\mathcal{O}(D_{\mathrm{diag}})$ has finite global
sections only when the covering degree $tN$ stays below the
discriminant threshold.

\paragraph{Orbit decomposition $(g_N, h_M)$.}
Slide lines~920--953 (Conjecture~1 expansion) assign each form to a
commuting K3-automorphism pair $(g_N, h_M)$ via the orbit
$X = (S \times E)/(\mathbb{Z}/N\mathbb{Z})$ with $g_N$ K3-symplectic of
order~$N$ and $h_M$ order~$M$:

\begin{center}
\small
\begin{tabular}{lll}
\hline
Form & $(g_N, h_M)$ orbit & Slide reference \\
\hline
$M_5(\Gamma_1,\nu_2) = \Delta_5$   & $(1, 1)$ trivial                    & line~912 Theorem ($N=1$) \\
$M_2(\Gamma_2,\nu_4)$              & $(g_2, 1)$                          & line~943 Case $N\in\{2,3,5,7\}$ \\
$M_3(\GammaT{1}{2},\nu_2)$         & $(1, h_2)$                          & line~943 \\
$M_1(\Gamma_3,\nu_6)$              & $(g_3, h_3)$                        & line~943 \\
$M_2(\GammaT{1}{3},\nu_2)$         & $(1, h_3)$                          & line~943 \\
$M_{1/2}(\Gamma_4,\nu_8)$          & $(g_4, h_4)$                        & lines~945--946 Case $N=4$ \\
$M_{3/2}(\GammaT{1}{4},\nu_4)$     & $(g_2, h_4)$ (mixed order)          & line~945 \\
$M_1(\GammaT{2}{2},\nu_4)$         & $(g_2, h_2)$ diagonal               & line~943 \\
\hline
\end{tabular}
\end{center}

\paragraph{Clery--Gritsenko 2013 theorem.}
\ClaimStatusTheorem{} Gritsenko, V.; Cl\'ery, F. ``Siegel modular forms
of genus~$2$ with the simplest divisor.'' Compos.\ Math.\ 149 (2013),
1563--1593, \S1 Theorem~1.1 (the eight-form classification): the
dd-Siegel forms across paramodular groups of genus~$2$ form a finite
set of eight; each is a Borcherds-additive lift of an explicit weak
Jacobi form of weight $0$ and index $t$ of the prescribed level and
character.

\paragraph{Mukai 1988 cross-reference.}
\ClaimStatusTheorem{} Mukai, S. ``Finite groups of automorphisms of K3
surfaces and the Mathieu group.'' Invent.\ Math.\ 94 (1988), 183--221,
classifies symplectic-automorphism orders on K3 into 11 classes
realised as subgroups of the Mathieu group $M_{23}$. The subset
$\{1, 2, 3, 4, 5, 6, 7, 8\}$ of orders~$\leq 8$ exactly matches the
eight $(g_N, h_M)$ pairs with $\mathrm{lcm}(N, M) \leq 8$ arising in
Gritsenko--Cl\'ery; the Mukai bound $|G| \leq 960$ and the
$M_{24}$-realisation (Kondo 1998) restrict to the Vol~III subset
$N \in \{1, 2, 3, 4, 6\}$ admitting full Borcherds-denominator control.

\paragraph{Programme $N \in \{1,2,3,4,6\}$ assignment.}
The subset of the eight forms for which Vol~III's
$\kBKM(\Phi_N) = c_N(0)/2$ identity is rigorously available comprises
five: $\Delta_5$ (weight~5, $N=1$, $\kBKM=5$),
$M_2(\Gamma_2,\nu_4)$ (weight~2, $\kBKM=2$),
$M_3(\GammaT{1}{2},\nu_2)$ (weight~3, $\kBKM=3$),
$M_2(\GammaT{1}{3},\nu_2)$ (weight~2, $\kBKM=2$),
$M_{3/2}(\GammaT{1}{4},\nu_4)$ (weight~$3/2$, $\kBKM = 3/2$). The
remaining three forms with $k \in \{1, 1/2\}$ ($M_1(\Gamma_3, \nu_6)$,
$M_{1/2}(\Gamma_4, \nu_8)$, $M_1(\GammaT{2}{2}, \nu_4)$) exit the
integer-weight regime and are conjectural within the programme.

\paragraph{References cited.}
Gritsenko, V.; Cl\'ery, F., Compos.\ Math.\ 149 (2013), 1563--1593;
Mukai, S., Invent.\ Math.\ 94 (1988), 183--221; Lorgat, R.
``A Borcherds lift of the weak Jacobi form $\phi_{0,1}$, generalized
Borcherds--Kac--Moody superalgebras and the Igusa cusp form $\Delta_5$,''
April 2020 preprint and accompanying presentation.
